\documentclass{article}

\usepackage[margin=0.75in,footskip=0.25in]{geometry}
\usepackage[T2A]{fontenc}
\usepackage[utf8]{inputenc}
\usepackage[serbianc, english]{babel}


\usepackage{titlesec}
\titlelabel{\thetitle.\quad}
\usepackage{array}
\usepackage{graphicx}
\usepackage{hyperref}
\usepackage{dirtytalk}
\newcommand{\jumpking}{\say{Jump King}~}

\title{\begin{otherlanguage*}{serbianc}{Испит из Основних оптимизационих алгоритама у инжењерству}\end{otherlanguage*}}
\date{\vspace{-5ex}}
% \date{}  % Toggle commenting to test

\begin{document}

\maketitle

\begin{otherlanguage*}{serbianc}
% Napomena
Напоменe: Испит траје до 150 минута. На усмену одбрану испитног задатка долази се са попуњеним овим документом и одговарајућим прилозима (фајловима). Студент попуњава: своје податке у табели, датум испита, као и садржаје одељака од 1 до 5, уносећи одговарајуће податке уместо објашњења за попуњавање. Задатак се ради коришћењем рачунара, а коришћено развојно окружење или одговарајући софтверски пакет морају бити доступни за време испита. Студент може понети и свој рачунар. Задатак носи до 30 поена.

% Tabela
\begin{center}
\begin{tabular}{ | m{2.5cm} | m{6cm}| m{2cm} | m{1.5cm} | m{1.5cm} | m{1.5cm} |}
\hline
Индекс\newline(година/број)& Презиме и име & Датум & Поени & Испит & Оцена\\
\hline
2021/0379 & Виктор Живковић & 08.08.2025. & & & \\[5ex]
\hline
\end{tabular}
\end{center}

\section{Опис проблема оптимизације}

Проблем који се решава је рачунање оптималних команди за прелазак игрице \jumpking
коришћењем генетичког алгоритма.
Идеја за овај проблем је потекла из јутјуб видеа са канала \say{Code Bullet}~\cite{codebulletyt},
у коме он програмира генетички алгоритам који игра игру и успева да победи.

\jumpking је платформска видео-игра у којој играч управља ликом(Краљем) који се
вертикално пење искључиво шетањем лево-десно и извођењем скокова чија се снага одређује дужином
држања тастера.
Недостатак контроле у ваздуху и могућност да и због најмање грешке играч изгуби велики део
напретка чине ову игру изазовном.

Тренутни светски рекорд у брзини прелажења игре у категорији за људе је 3 минута, 59 секунди и
598 стотинки \cite{speedruncom}, а у категорији \say{TAS}(Tool-Assisted Speedrun) је
3 минута, 43 секунде и 973 стотинке \cite{tasrecord}.

Како је питање о низу уноса команди, које могу бити један од три тастера на тастатури, овај проблем
се сврстава у групу SAT проблема.

Код оригиналне игре није доступан, тако да сам се ја одлучио за open-source рекреацију написану
у програмском језику Python~\cite{jkathome}.


\section{Дефиниција оптимизационе функције и оптимизационог простора}
Систем управљања сведен је на само три дугмета - лево за шетање лево, десно за шетање десно
и горе за скок, јачина скока одређена је дужином држања тастера за скок.

Ради дефинисања оптимизационих променљивих и оптимизационог простора,
ток играња једне \say{партије} је подељен на акције.
Једна акција састоји се из
команде (лево, десно, скок улево, скок удесно) и њеног трајања.
Акција је представљена једним целим бројем
од 8 бита, где 2 бита највеће тежине представљају код команде, а осталих 6 бита
представљају дужину трајања акције, у \say{фрејмовима}(генерисаним сликама).

Једна партија састоји се из низа акција, којих може бити произвољан број.
Оптимизациони простор дефинисан је као низ акција, посматраних једна за другом,
као низ од 8 * N битова, где је N број акција.
Како димензионалност оптимизационог простора расте линеарно са бројем акција,
а комплексност проблема експоненцијално зависи од димензионалности,
партија је подељена на делове од 5 акција које се независно посматрају.

Ова одлука има смисла са обзиром да је циљ игре увек попети се што више за
што краћи временски период.
Међутим, постоји проблем са нивоима који захтевају привремено спуштање да
би се ниво прешао, у оваквим нивоима постоји шанса да алгоритам не нађе
оптимално решење ако део који захтева спуштање не буде потпуно обухваћен
једним оптимизационим блоком од 5 акција.
Конкретно у овом раду алгоритам још увек није стигао до таквог нивоа тако
да је ово проблем који се може јавити у будућности развијања.

За рачунање оптимизационе функције неопходна је симулаицја која обухвата
све акције из блока, из које само
последња позиција играча и укупно време у игри улазе у прорачун
оптимизационе функције.

Постоје 2 критеријума, висина коју је играч достигао и време које му је било
потребно за то.
Критеријуми се комбинују у једну оптимизациону функцију, тако да се критеријуму
висине даје много већи значај. Један ниво је висине једнe слике(екрана), што је са
нормалним подешавањима 720 пиксела.
Време се мери у секундама и потом множи са 0,1 што значи да на крају игре ако
је решење оптимално или близу оптималном биће отприлике 2,7 што је значајно
мање од првог критеријума.
Оптимизациона функција се рачуна на начин описан једначином~\ref{eq:f},
где је:
\begin{itemize}
    \item $l$  тренутни ниво
    \item $h$  висина екрана у пикселима
    \item $y$  позиција играча на екрану у пикселима по вертикалној оси са почетком на врху екрана
    \item $t$  време.
\end{itemize}

\begin{equation}
    \label{eq:f}
    f = (l+1) * h - y - 0.1 * t
\end{equation}



\section{Оптимизациони алгоритам или алгоритми}
За решавање проблема коришћен је генетички алгоритам представљен на предавањима~\cite{predavanja}.
Величина популације $N_{pop}$ је 300, а број генерација је 150.

Турнир селекцијом са величином турнира 1 селектује се $N_{pop}/{5}$ родитеља.

Насумично се бирају 2 родитеља који се онда укрштају са
вероватноћом 0,8.
Имплементирано је бинарно укрштање са једном тачком пресека, за сваки пар
родитеља укрштање се ради 2 пута, па пар родитеља производи четворо деце.

Вероватноћа мутације сваког бита у генерацији је 0,1.

За рачунање оптимизационе функције потребно је симулирање целе партије
са потезима једне јединке, што је захтевно.
Ради убрзавања овог процеса искључена је
генерација слика приликом симулације играња и
имплементирана је паралелизација која омогућава покретање
задатог броја симулација у исто време. Овај број зависи од рачунара
на коме се покреће оптимизација, за мој кућни рачунар оптималан број
конкурентних симулација је 6.

Почетно решење генерисано је потпуно случајно.


\section{Резултати}
Могуће је погледати репризу игре коју је одиграо алгоритам, покретањем фајла
\texttt{Brain.py}.
Алгоритам успешно прелази првa 2 нивоа, међутим на 3. нивоу се зауставља.
Долази до дела где је потребно направити за редом 3 веома прецизна скока, која покрећу карактера искључиво
хоризонтално да би после тога било могуће направити вертикални скок са којим ће карактер заправо напредовати
(у смислу оптимизационе функције).
У случају грешке приликом било којег од ових \say{хоризонталних} скокова карактер пада на почетак прошлог
нивоа, што значи да у овом делу оптимизационог простора постоји веома узак минимум окружен униформним вредностима
опт. функције. Из тог разлога није могуће да алгоритам закључи где је минимум осим ако га
случајно не погоди путем мутација, међутим како је овај минимум веома узак шансе за то су веома мале.

Прво од могућих решења овог проблема је повећавање величине популације, самим тим постоји више јединки и веће
су шансе да ће нека од њих случајно упасти у овај минимум. Међутим, проблем са тим је што комплексност веома брзо расте
и алгоритам почиње да захтева превише времена за извршавање.

Следећа идеја за решење је додавање неког регуларизационог сабирка у оптимизациону функцију који би требао
да \cite{прошири} овај минимум. Комплексност ове идје је у смишљању ове функције.


% Време потребно за оптимизацију до овог тренутка било је:

% Дакле рачунајући да постоји     нивоа, укупно естимирано време оптимизације би било:


\section{Прилози}

Код се налази на
\href{https://github.com/Sslowepoke/Jumpkingathome}%
{Github линку \texttt{https://github.com/Sslowepoke/Jumpkingathome}}
и у фолдеру \texttt{JumpKingAtHome/}.

Једини фајлови који су мењани у односу на оригинални код игрице су
\texttt{Brain.py} и \texttt{JumpKing.py}.
У фајлу \texttt{Brain.py} налази се имплементација оптимизационог алгоритма, док се
у фајлу \texttt{JumpKing.py} налази модификован интерфејс са остатком игре.

У случају покретања фајла \texttt{Brain.py} врши се оптимизација,
а у случају покретања \texttt{JumpKing.py} могуће је играти игру нормално.

Свака етапа оптимизације чува резултате у фајлу чије је име време завршавања оптимизације
а налази се у фолдеру \texttt{Saves/}, док се финално решење састојано из надовезаних резултата
етапа оптимизације налази у фајлу \texttt{Saves/latest.txt}.



\end{otherlanguage*}


\selectlanguage{serbianc}
\begin{thebibliography}{9}

\bibitem{codebulletyt}
    Code Bullet,
     AI Learns to play JUMP KING,
     YouTube. \url{https://youtu.be/DmQ4Dqxs0HI}

\bibitem{speedruncom}
    Speedrun.com
    \url{https://www.speedrun.com/jumpking}

\bibitem{tasrecord}
    \url{https://www.youtube.com/watch?v=LZlo6TzL7N4}

\bibitem{jkathome}
    senweim,
    JumpKingAtHome,
    \url{https://github.com/senweim/JumpKingAtHome}

\bibitem{predavanja}
    Драган Олћан,
    Предавања из предмета ООА,
    2025.






\end{thebibliography}

\end{document}
